\section{Síntesis y ampliación}

La semana 2 profundizó en el núcleo técnico del Coding Agent. Desde construir manualmente un Coding Agent mínimo hasta entender cómo el protocolo MCP estandariza la integración de herramientas, el curso mostró la implementación de ingeniería detrás de la «magia» de las herramientas de programación con IA.

\subsection{Tabla de síntesis}

\begin{table}[H]
\centering
\begin{tabular}{p{2.8cm}p{4.2cm}p{4.0cm}}
\toprule
\textbf{Tema} & \textbf{Conclusión central de la semana 2} & \textbf{Implicación de ingeniería} \\
\midrule
Agent Loop mínimo & LLM + tool call + observation ya bastan para formar un ciclo básico & El principio es simple, pero el contexto y el manejo de errores determinan la usabilidad \\
Experiencia de Claude Code & El reminder, la precarga de contexto y la distribución con subagent son fundamentales & El diseño del producto incide directamente en la estabilidad del modelo \\
Protocolo MCP & La integración de herramientas pasa de $M \times N$ a $M + N$ & La innovación del ecosistema puede pasar de integraciones privadas a interfaces públicas \\
Herramientas AI-native & Una vez unificado el protocolo, el diseño de las herramientas mismas se vuelve el cuello de botella & La API necesita ser más plana, recuperable e interpretable \\
\bottomrule
\end{tabular}
\caption{La semana 2 se parece más a una clase de diseño de infraestructura para Coding Agent que a una simple clase de uso de herramientas}
\end{table}

\subsection{Puntos clave}

\begin{itemize}
    \item Coding Agent = ciclo de LLM + Tool Calling; el núcleo es simple, pero los detalles de ingeniería determinan la calidad
    \item Diseños clave de Claude Code: precarga de contexto, system reminders para evitar el drift, distribución con subagent
    \item MCP reduce la integración de $M \times N$ a $M + N$, usando JSON-RPC para estandarizar la descripción de herramientas
    \item Arquitectura de MCP: Host $\rightarrow$ Client $\rightarrow$ Server $\rightarrow$ Tool
    \item Límite actual: la capacidad del agente para manejar múltiples herramientas es limitada; la API necesita un diseño AI-native
\end{itemize}

\subsection{Lecturas adicionales}

\begin{itemize}
    \item Documentación oficial de MCP: \url{https://modelcontextprotocol.io}
    \item Blog de lanzamiento de MCP de Anthropic: \url{https://www.anthropic.com/news/model-context-protocol}
    \item Especificación de JSON-RPC 2.0: \url{https://www.jsonrpc.org/specification}
    \item Language Server Protocol: \url{https://microsoft.github.io/language-server-protocol/}
\end{itemize}

%% --- Fin del contenido --- %%

\end{document}
